\documentclass[10pt]{article}
\usepackage{icmlko}

\icmlmeta{IP 파운데이션 월드모델}{181}{축이 아니라 라벨이 달랐다}{노트 180이 ``위키 축의 부호가 도메인마다 반대''를 남기고 원작 가설을 걸었다. 반증됐다 --- 속편이었고, 속편을 벌하는 것은 축이 아니라 라벨이었다}

\begin{document}

\twocolumn[
\icmlheader

\begin{abstract}
노트 180이 ``애니 · 웹툰 · 만화는 한국어 위키에서 \emph{원작} 문서에 붙어서
음수일 것''을 걸고 ``분류로 검정할 수 있다''로 닫았다. 검정했다.
\textbf{반증됐다} --- 애니에서 원작형 문서 46건이 $\rho{=}-0.58$ 이고
자신형 64건이 $-0.41$ 로 \textbf{둘 다 음수}다. 문서 종류는 크기를 바꾸지
부호를 안 바꾼다. 다른 갈래를 댔다. \textbf{속편이다.} 애니 속편 240건이
$\rho{=}\mathbf{-0.28}$ 이고 신작 442건이 $\mathbf{+0.04}$ 다 ---
\textbf{음수가 속편에만 있다.} 까닭이 둘로 갈린다: 속편은 사전 위키
조회수가 훨씬 크고($\log$ 수준 2.28 대 0.00 · $p{=}2.7\times10^{-35}$)
한국 라벨은 더 낮다(2.14 대 2.34 · $p{=}1.1\times10^{-6}$). 그리고
\textbf{세계애니가 대조군이 된다} --- 거기서도 속편이 훨씬 유명한 것은
같은데(6.56 대 5.35 · $p{=}3.9\times10^{-24}$) \textbf{라벨이 안 낮다}
(4.77 대 4.74 · $p{=}1.00$). 그래서 속편 $\rho$ 가 $+$0.59 로 오히려 더
크다. \textbf{두 도메인의 차이는 축이 아니라 라벨에 있다} --- 애니 라벨
(한국 시청)은 22기를 벌하고 세계애니 라벨은 안 벌한다. 그러면 속편 여부를
축으로 세우면 되나. 세웠다 --- 덮음 100\% · 수집 비용 0 · 누수 0 인 깨끗한
축이다. \textbf{$+$0.0022($t{=}0.95$)로 안 갈린다.} 속편 정보는 이미 위키
조회수 안에 들어 있었다 --- \textbf{노트 175에서 ``문서가 있었나''를 명시
축으로 세웠을 때 $+$0.0000 이었던 것에 이어 세 번째다.} 마지막으로 남는
것: \textbf{웹툰과 만화의 음수는 속편이 아니다} --- 속편 몫이 0.4\%와
1.2\%다. 설명 못 한 채로 둔다.
\end{abstract}
]

\begin{figure*}[t]
\centering
\includegraphics[width=\textwidth]{figs/label.pdf}
\caption{왼쪽 --- 애니만 속편에서 뒤집힌다. 가운데 --- 원작 가설은
반증됐다. 오른쪽 --- 명시 축은 잉여다.}
\end{figure*}

\section{원작 가설은 반증됐다}

\begin{table}[t]
\centering\tsize
\caption{붙은 문서의 종류별 위키 수준 $\leftrightarrow$ 라벨.}
\begin{tabular}{lrrrr}
\toprule
& \multicolumn{2}{c}{애니} & \multicolumn{2}{c}{세계애니} \\
& n & $\rho$ & n & $\rho$ \\
\midrule
원작형 & 46 & \textbf{$-$.58} & 0 & --- \\
자신형 & 64 & \textbf{$-$.41} & 45 & $+$.14 \\
둘 다 & 330 & $-$.11 & 1232 & $+$.28 \\
\bottomrule
\end{tabular}
\end{table}

\claimx{\textbf{애니는 셋 다 음수다.} 원작 문서에 붙었든 그 작품 자신의
문서에 붙었든 부호가 같다 --- \textbf{문서 종류는 크기를 바꾸지 부호를
안 바꾼다.}}{노트 180이 걸어 둔 ``반대 가능성''이 그대로 나왔다 ---
``부호가 문서 종류가 아니라 도메인 자체에서 올 수 있다.'' \textbf{내가 건
가설 둘 중 안 유리한 쪽이 맞았다.}}

\section{속편이었다}

\begin{table}[t]
\centering\tsize
\caption{속편 대 신작. 옳은 매칭만.}
\begin{tabular}{lrrrr}
\toprule
& \multicolumn{2}{c}{애니} & \multicolumn{2}{c}{세계애니} \\
& n & $\rho$ & n & $\rho$ \\
\midrule
\textbf{속편} & 240 & \textbf{$-$.281} & 223 & \textbf{$+$.586} \\
\textbf{신작} & 442 & \textbf{$+$.045} & 1332 & $+$.179 \\
\midrule
속편 몫 & \multicolumn{2}{c}{46.6\%} & \multicolumn{2}{c}{13.2\%} \\
\bottomrule
\end{tabular}
\end{table}

\claimx{\textbf{애니의 음수가 속편에만 있다.} 신작 442건은 $+$0.045로
사실상 0 이다. 노트 180이 잰 애니 전체 $-$0.18 은 속편 46.6\%가 만든
것이다.}{노트 149가 프랜차이즈 되돌리기를 넣으면서 ``원작 인기를 업고
나오는 2기와 맨몸으로 나오는 신작을 가르는 축이 된다''고 적었는데,
\textbf{정말로 그 둘이 갈렸다} --- 다만 같은 방향이 아니라 반대
방향으로.}

\section{다른 것은 라벨이었다}

\begin{table}[t]
\centering\tsize
\caption{속편이 신작과 어떻게 다른가. 중앙값.}
\begin{tabular}{lrrl}
\toprule
& 속편 & 신작 & \\
\midrule
\multicolumn{4}{l}{\emph{애니}} \\
$\log$ 위키 수준 & 2.28 & 0.00 & $p{=}2.7{\times}10^{-35}$ \\
\textbf{라벨} & \textbf{2.14} & \textbf{2.34} & $p{=}1.1{\times}10^{-6}$ \\
\multicolumn{4}{l}{\emph{세계애니}} \\
$\log$ 위키 수준 & 6.56 & 5.35 & $p{=}3.9{\times}10^{-24}$ \\
\textbf{라벨} & \textbf{4.77} & \textbf{4.74} & \textbf{$p{=}1.00$} \\
\bottomrule
\end{tabular}
\end{table}

\claimx{\textbf{축은 두 도메인에서 똑같이 행동한다} --- 속편이 훨씬 유명하다
($p{<}10^{-23}$, 양쪽 다). \textbf{다른 것은 라벨이다} --- 애니 라벨은
속편을 벌하고($p{=}1.1\times10^{-6}$) 세계애니 라벨은 안 벌한다
($p{=}1.00$).}{읽기가 하나 있다 --- 애니 라벨은 \emph{그 분기의} 한국
시청이고, 세계애니 라벨은 작품에 누적된 값이다. ``명탐정 코난 22기''는
프랜차이즈가 크지만 그 분기를 한국에서 보는 사람은 적다. \textbf{누적
라벨은 그 벌을 안 받는다.}}

\claimx{\textbf{이건 자료 잘못이 아니라 세계 사실이다.} 노트 175 · 164가
``축 이름이 내용과 다르다''를 잡은 것과 다르다 --- 여기서는 축이 옳게
재고 있고 \textbf{두 도메인의 라벨이 다른 것을 묻고 있다.}}{그래서 도메인
사이에 계수를 나눠 쓰는 것이 옳은지가 흔들린다. 풀링 자체가 ``라벨이
같은 것을 묻는다''를 가정하는데 여기서는 안 그렇다.}

\section{그런데 축으로 세우면 잉여다}

\begin{table}[t]
\centering\tsize
\caption{\texttt{is\_sequel} --- 덮음 100\% · 수집 비용 0 · 누수 0.}
\begin{tabular}{lrrl}
\toprule
& 차 & sd & \\
\midrule
F18 배깅 & $+$.0022 & .0023 & $t{=}0.95$ \\
F6 능형 & $-$.0016 & .0023 & $t{=}-0.72$ \\
\midrule
애니 & $+$.0078 & & 속편 46.6\% \\
팝업 & $-$.0097 & & 속편 없음 \\
\bottomrule
\end{tabular}
\end{table}

\claimx{\textbf{안 갈린다.} 예측한 도메인(애니 $+$0.0078)은 맞았는데 판이
안 움직인다 --- 속편 정보가 이미 위키 조회수 안에 있다(속편은 프랜차이즈
문서에 붙어 조회수가 크다).}{\textbf{세 번째다} --- 노트 175에서 ``문서가
있었나''를 명시 축으로 세웠을 때 $+$0.0000($t{=}0.0$)이었고, 노트 166에서
제작사 이력을 복제했을 때 공선성만 얻었다. \textbf{모형이 이미 쓰고 있는
것을 명시 축으로 세우면 잉여다.}}

\claimx{내 예측이 반쯤 맞았다 --- ``F18 $+$0.003$\sim$$+$0.010''은 넘겼고
``F6 $+$0.000$\sim$$+$0.003''도 넘겼지만, 걸어 둔 반대 가능성(``이미 위키
수준 안에 들어 있다'')이 실제로 일어났다.}{\textbf{반대 가능성을 미리 적어
두는 것이 값을 한다} --- 노트 180에서도 걸어 둔 쪽이 맞았다. 두 사이클
연속이다.}

\begin{table}[t]
\centering\tsize
\caption{남은 것.}
\begin{tabular}{ll}
\toprule
\textbf{웹툰 · 만화} & 속편이 0.4\% · 1.2\% --- 음수가 미설명 \\
풀링 & 라벨이 다른 것을 묻는데 계수를 나눠 쓰나 \\
상호작용 & 속편$\times$위키 수준을 곱해 넣으면 \\
\bottomrule
\end{tabular}
\end{table}

두 번째 줄이 제일 크다. 이 판의 얼개는 열 도메인의 계수를 나눠 쓰는
풀링인데(노트 5), 애니 라벨은 분기별 시청이고 세계애니 라벨은 누적
회원이다. \textbf{같은 축이 두 라벨에 반대로 붙는 것을 봤으니, 풀링이
가정하는 ``라벨이 같은 것을 묻는다''가 적어도 한 쌍에서는 거짓이다.}
노트 141이 라벨 상관 0.70 을 누수 문턱으로 쓴 자리와 반대편이다 --- 거기서는
너무 닮아서 문제였고 여기서는 너무 달라서 문제다.

\end{document}
